% =====================================================================
%  Master's Thesis Section: Self-Written SLP Solver
%  --------------------------------------------------------------------
%  Drop-in section for the "Optimization" chapter. Compile as a
%  standalone document or \input{} the \section blocks into a chapter.
%  --------------------------------------------------------------------
%  Required preamble packages (already in most thesis templates):
%      amsmath, amssymb, booktabs, enumitem, hyperref
% =====================================================================

\documentclass[11pt,a4paper]{article}

\usepackage[margin=25mm]{geometry}
\usepackage{amsmath}
\usepackage{amssymb}
\usepackage{booktabs}
\usepackage{enumitem}
\usepackage{hyperref}

\title{Self-Written Sequential Linear Programming Solver:\\
       A Theoretical Background for the Master's Thesis}
\author{DMS4 Optimization Chapter}
\date{\today}

\begin{document}

\maketitle

% =====================================================================
\section{Self-Written Sequential Linear Programming Solver}
\label{sec:slp_solver}

\subsection{Motivation and Project Context}

In this thesis the structural mass of a wind-turbine support structure is
minimized subject to load-, deflection-, and buckling-related inequality
constraints. The design vector $x \in \mathbb{R}^n$ contains geometrical
parameters with very different physical units --- diameters in the order of
hundreds of millimetres next to wall thicknesses in the order of a few
millimetres. The objective and constraints are evaluated by an external finite
element model, so the solver only sees them as black boxes for which gradients
must be approximated by finite differences.

Off-the-shelf optimizers such as \textsc{PySLSQP} or \textsc{SciPy}'s
\texttt{minimize} solve this class of problem efficiently, but they are
opaque: when a structural step is rejected or a solution is locally
infeasible, it is difficult to tell why. To gain a deeper understanding of
the underlying mechanics --- and to be able to discuss them in this thesis ---
a transparent Sequential Linear Programming (SLP) solver was implemented from
scratch. The implementation deliberately keeps the algorithmic ingredients
modular and observable: every iteration logs the predicted-versus-actual
quality ratio, the linearization slacks, the active move limits, and the
current penalty parameter. This makes the solver well suited for didactic use
and ideal for the thesis discussion.

The methods used in the solver are not new; they are textbook ideas
\citep{MartinsNing2021,NocedalWright2006,Fletcher1987}. The contribution of
this section is to present those ideas in a single, coherent framework and to
explain \emph{why each ingredient is needed} for a robust SLP method to work
on structural design problems.

\subsection{Problem Formulation}

The optimization problem solved in this work is written in the
\textsc{PySLSQP} convention with non-negative inequality constraints:
\begin{align}
    \min_{x \in \mathbb{R}^n} \quad
        & f(x), \\
    \text{subject to} \quad
        & c_i(x) \geq 0, \qquad i = 1, \dots, m, \\
        & x^L \leq x \leq x^U,
\end{align}
where $f : \mathbb{R}^n \to \mathbb{R}$ is the structural mass and
$c : \mathbb{R}^n \to \mathbb{R}^m$ is the vector of nonlinear inequality
constraints. A constraint is feasible when its value is non-negative and
violated otherwise. The bounds $x^L$ and $x^U$ enforce the physical limits of
the geometric design variables.

\subsection{Method Overview}

Sequential Linear Programming approximates the nonlinear problem by a
\emph{sequence of linear subproblems}. At iteration $k$ the objective and
constraints are linearized around the current design $x_k$:
\begin{align}
    f(x_k + \Delta x) &\approx f(x_k) + g(x_k)^{\!\top} \Delta x, \\
    c(x_k + \Delta x) &\approx c(x_k) + J(x_k) \Delta x,
\end{align}
where $g = \nabla f$ and $J$ is the constraint Jacobian. A linear program
(LP) then proposes the next step $\Delta x$. Because a linear model is only
locally accurate, every step must be globalized: it must be \emph{accepted}
or \emph{rejected} based on what the true nonlinear functions do at the
trial point. The implemented solver accomplishes this through six
ingredients that we discuss in detail in the following subsections:

\begin{enumerate}[nosep]
    \item \textbf{Variable scaling} --- equalize the influence of variables
          with very different units on the LP.
    \item \textbf{Slack-relaxed LP subproblem} --- guarantee that the LP
          always has a solution even when the linearized constraints cannot
          be satisfied within the current move limits.
    \item \textbf{$L_1$ merit function} --- combine objective and
          constraint violation into a single scalar criterion.
    \item \textbf{Inexact (backtracking) line search} --- accept only
          trial steps that produce sufficient merit decrease.
    \item \textbf{Adaptive move limits} --- adjust the local trust region
          based on the agreement between linear model and nonlinear
          response.
    \item \textbf{Adaptive penalty update} --- increase the importance of
          feasibility when constraints remain active.
\end{enumerate}

\subsection{Variable Scaling}
\label{sec:scaling}

\subsubsection*{Theory}

In structural design the components of $x$ are physical lengths whose
magnitudes can differ by several orders. A radius may be of order
$200~\text{mm}$ while a wall thickness may be of order $2~\text{mm}$. The
gradients of $f$ and $c$ inherit those magnitudes through the chain rule.
If the LP is solved in the physical variables, a single ill-scaled
variable can dominate the search direction purely because of its unit, not
because it actually drives the design. \citet[Sec.~4.4]{MartinsNing2021}
emphasize that variable scaling is one of the most consequential numerical
practices in gradient-based design optimization, and \citet[Sec.~6.5]{NocedalWright2006}
note that all gradient-based methods are sensitive to scaling because the
geometry of the feasible set in scaled space drives the conditioning of the
subproblem.

The implemented solver uses a bound-based linear scaling that maps every
bounded variable into the unit interval:
\begin{align}
    y_j = \frac{x_j - x^L_j}{x^U_j - x^L_j}, \qquad
    x_j = x^L_j + (x^U_j - x^L_j)\, y_j.
\end{align}
For variables without finite bounds a fallback scale of
$\max(1, |x_j|)$ is used. In compact notation,
$x = x^{\mathrm{off}} + D\, y$, where $D$ is the diagonal scale matrix.
By the chain rule, gradients in the scaled space are
\begin{align}
    g_y = D\, g_x, \qquad J_y = J_x\, D.
\end{align}
The LP step is then computed in $y$-space and converted back to physical
units before any nonlinear evaluation.

\subsubsection*{Why this matters in this project}

\begin{itemize}[nosep,leftmargin=1.5em]
    \item \textbf{Conditioning.} In scaled coordinates each design
          variable has unit range; the LP no longer favours large
          variables simply because they have a large coefficient in $g$
          or $J$.
    \item \textbf{Move limits become physically meaningful.} A move
          limit of $\delta = 0.1$ in $y$-space corresponds to ten percent
          of each variable's allowable physical range, regardless of unit.
          Without scaling, a single move limit cannot be appropriate for
          radii and thicknesses simultaneously.
    \item \textbf{Tolerances are unit-independent.} Convergence checks
          based on $\|g_y\|_\infty$ are invariant under a change of
          physical units, which is essential for a structural model
          whose variables come from different design domains.
\end{itemize}

\subsection{Physical vs.\ Scaled Space --- Where Each Quantity Lives}
\label{sec:phys_vs_scaled}

The solver deliberately keeps two coordinate systems active at the same
time: the LP subproblem and its move limits live in the scaled
$y$-space, while the merit function, the line-search acceptance test,
the trust-region ratio $\rho$, and the convergence tests are evaluated
in the physical $x$-space. The bridge between the two is the chain
rule $g_y = D g_x$ and $J_y = J_x D$, which guarantees that the
linearized changes
\begin{align}
    g_y^{\!\top}\Delta y \;=\; g_x^{\!\top}\Delta x,
    \qquad
    J_y\,\Delta y \;=\; J_x\,\Delta x,
\end{align}
come out in physical units even though $\Delta y$ is scaled. The
following table summarizes which quantity lives in which space and is
referred to throughout the rest of this section.

\begin{table}[h]
\centering
\small
\begin{tabular}{p{0.30\textwidth} p{0.32\textwidth} l}
\toprule
Quantity & Symbol & Space \\
\midrule
Design variable                       & $x$                                & physical \\
Scaled design variable                & $y = D^{-1}(x - x^{\mathrm{off}})$ & scaled \\
Bounds                                & $x^L, x^U$ \;/\; $y^L, y^U$        & physical / scaled \\
Move limit (trust region)             & $\delta$                           & scaled \\
LP step direction                     & $\Delta y$                         & scaled \\
Equivalent physical step              & $\Delta x = D\,\Delta y$           & physical \\
Slack variables                       & $s$                                & physical (constraint units) \\
Penalty                               & $M$                                & $[f\text{-units}]/[c\text{-units}]$ \\
Objective                             & $f(x_k)$                           & physical \\
Constraint vector                     & $c(x_k)$                           & physical \\
Physical gradients                    & $g_x = \nabla_x f,\; J_x = \nabla_x c$ & physical \\
Scaled gradients                      & $g_y = D g_x,\; J_y = J_x D$       & scaled \\
LP cost coefficient on $\Delta y$     & $g_y$                              & scaled gradient \\
LP RHS of linearized constraint       & $c_k$                              & physical \\
Linearized $f$-change                 & $g_y^{\!\top}\Delta y$             & physical (by chain rule) \\
Linearized $c$-change                 & $J_y\,\Delta y$                    & physical (by chain rule) \\
Merit function $\Phi(x;M)$            & $f + M\!\sum\!\max(0,-c_i)$        & physical \\
Predicted / actual merit decrease     & $\Delta\Phi_{\mathrm{pred}}, \Delta\Phi_{\mathrm{act}}$ & physical \\
Trust-region ratio                    & $\rho = \Delta\Phi_{\mathrm{act}}/\Delta\Phi_{\mathrm{pred}}$ & physical \\
Step-norm convergence test            & $\|x_{k+1}-x_k\|_2$                & physical \\
Gradient-norm convergence test        & $\|g_y\|_\infty$                   & scaled \\
\bottomrule
\end{tabular}
\caption{Mapping of every algorithmic quantity to physical or scaled
space. The LP is solved in scaled coordinates for conditioning; the
nonlinear acceptance machinery is evaluated in physical coordinates so
that decisions reflect the true engineering problem.}
\label{tab:phys_vs_scaled}
\end{table}

A useful mnemonic is that everything related to \emph{building the
LP step} (variables, gradients, move limits, bounds) is in
$y$-space, while everything related to \emph{judging the step}
(merit, $\rho$, feasibility check, step norm) is in $x$-space.

\subsection{The Linearized Subproblem}

At iteration $k$, write the trust-region (move-limit) constraint in
scaled coordinates as $\Delta y \in [-\delta, \delta]$ and combine it with
the global bounds $y^L \leq y_k + \Delta y \leq y^U$. The decision
variables of the LP, the move limits, and the gradient $g_{y,k}$ are
all in \emph{scaled} space, but the right-hand side $c_k$ of the
linearized inequality and the LP objective value
$g_{y,k}^{\!\top}\Delta y$ are in physical units --- the chain-rule
relation $g_y = D g_x$ guarantees that
$g_{y,k}^{\!\top}\Delta y = g_{x,k}^{\!\top}\Delta x$, so the LP
minimizes the same physical $f$-decrease as a hypothetical LP solved
directly over $\Delta x$. A first attempt at the SLP subproblem is
the pure LP
\begin{align}
    \min_{\Delta y} \quad
        & g_{y,k}^{\!\top} \Delta y, \\
    \text{subject to} \quad
        & c_k + J_{y,k}\, \Delta y \geq 0, \\
        & \max(y^L - y_k,\, -\delta) \leq \Delta y \leq
          \min(y^U - y_k,\, \delta).
\end{align}
This subproblem has a major weakness: as soon as the linearized
constraints $c_k + J_{y,k}\,\Delta y \geq 0$ cannot be satisfied within
the move limits, the LP becomes infeasible and the solver stalls. This is
not a hypothetical situation in structural design --- it occurs as soon as
the current iterate is infeasible, which is normal in early iterations.
The remedy is the introduction of slack variables.

\subsection{Slack Variables and the Slack-Relaxed LP Subproblem}
\label{sec:slacks}

\subsubsection*{Theory}

A slack variable is a non-negative auxiliary variable added to a constraint
to absorb the amount by which the constraint cannot be satisfied. Slacks
have a long history in linear programming and in the elastic-mode and
$\ell_1$-penalty SQP literature
\citep[Ch.~17--18]{NocedalWright2006},\citep{Fletcher1987}. The classical
$\ell_1$ formulation introduces $s \geq 0$ such that
\begin{align}
    c_k + J_{y,k}\,\Delta y + s \geq 0,
\end{align}
and adds a penalty term $M \sum_i s_i$ to the LP objective with $M > 0$
large. The resulting slack-relaxed LP solved by the implementation is
\begin{align}
    \min_{\Delta y, \, s} \quad
        & g_{y,k}^{\!\top} \Delta y + M \sum_{i=1}^{m} s_i, \\
    \text{subject to} \quad
        & c_k + J_{y,k}\,\Delta y + s \geq 0, \\
        & s_i \geq 0, \qquad i = 1, \dots, m, \\
        & \max(y^L - y_k,\, -\delta) \leq \Delta y \leq
          \min(y^U - y_k,\, \delta).
\end{align}
The standard upper-bound form used by linear-programming solvers is
recovered by writing the linearized constraint as
$-J_{y,k}\,\Delta y - s \leq c_k$, which is exactly the inequality
implemented in the code. Note again the mixed bookkeeping: the LP
decision variables $\Delta y$ and the scaled Jacobian $J_{y,k}$ live
in $y$-space, while the right-hand side $c_k$ and the slack vector
$s$ are in physical constraint units (so that $s_i$ literally
absorbs as many physical units of violation as required to make the
linearized residual non-negative).

\subsubsection*{Why slacks are good}

\begin{itemize}[nosep,leftmargin=1.5em]
    \item \textbf{Always feasible LP.} The relaxed LP has at least the
          trivial point $\Delta y = 0$, $s_i = \max(0, -c_{k,i})$ in its
          feasible set. The optimizer therefore never breaks down because
          of a momentarily infeasible iterate. This is essential in
          structural problems where the starting design is rarely strictly
          feasible.
    \item \textbf{Slacks reveal information.} A large optimal slack
          $s_i^\star$ is a quantitative signal that constraint $i$ cannot
          be satisfied locally within the present move limits. The solver
          uses $\sum s_i$ and $\max s_i$ as diagnostics: persistently
          large slacks indicate that either the move limits should
          shrink, the penalty should grow, or the linearization is poor
          (see Sections~\ref{sec:adaptive_move}, \ref{sec:adaptive_penalty}).
    \item \textbf{Penalty discourages artificial feasibility.} The term
          $M \sum s_i$ in the LP objective makes the optimizer prefer
          steps that satisfy the linearized constraints \emph{without}
          using slack whenever possible, while still allowing slack to
          be used when the alternative is no step at all. This is the
          $\ell_1$-penalty principle of \citet{Fletcher1987} adapted to
          the LP setting.
    \item \textbf{Connection to the merit function.} The same penalty
          $M$ appears in the nonlinear merit function defined in
          Section~\ref{sec:merit}. This consistency means that the
          subproblem and the acceptance test agree on how much
          infeasibility is tolerable.
\end{itemize}

\subsection{The $L_1$ Merit Function}
\label{sec:merit}

To decide whether a candidate step from the LP should be taken, the
nonlinear objective and the nonlinear constraint violation are combined
into a single scalar:
\begin{align}
    \Phi(x; M) \;=\; f(x) \;+\; M \sum_{i=1}^{m} \max\bigl(0, -c_i(x)\bigr).
\end{align}
This is the classical $\ell_1$ exact penalty function
\citep[Sec.~17.3]{NocedalWright2006}. \citet[Sec.~5.5]{MartinsNing2021}
discuss the merit function as the standard mechanism for transforming a
constrained acceptance test into an unconstrained one in SQP/SLP
globalization. The merit function and every quantity derived from it
($\Delta\Phi_{\mathrm{pred}}$, $\Delta\Phi_{\mathrm{act}}$, $\rho$,
$V$) is evaluated in \emph{physical} space: it takes physical $f(x)$
and physical $c(x)$ as inputs and returns a value with physical
$f$-units. The penalty $M$ therefore carries units of
$[f\text{-units}]/[c\text{-units}]$ and is identical to the LP slack
penalty introduced in Section~\ref{sec:slacks}, which keeps the LP
and the line search consistent.

\subsubsection*{Why the merit function matters}

\begin{itemize}[nosep,leftmargin=1.5em]
    \item It produces a single number on which a line search can act.
    \item It allows trade-offs: an objective decrease that worsens
          feasibility is only accepted when it does not worsen feasibility
          too much, weighted by $M$.
    \item It is exact for sufficiently large $M$: in particular, if $M$
          exceeds the largest absolute Lagrange multiplier of the
          original problem, then a local minimizer of the merit function
          coincides with a local minimizer of the constrained problem
          \citep[Thm.~17.3]{NocedalWright2006}. This justifies the
          adaptive penalty strategy in Section~\ref{sec:adaptive_penalty}:
          we may not know the multipliers, but we can grow $M$ until the
          merit function inherits this property.
\end{itemize}

\subsection{Inexact (Backtracking) Line Search}
\label{sec:line_search}

\subsubsection*{Theory}

A line search replaces the full LP step by a fraction of it. Given the
direction $\Delta y$ from the LP, the trial point is constructed in
\emph{scaled} coordinates and then mapped back to \emph{physical}
coordinates before any nonlinear evaluation:
\begin{align}
    y_{\mathrm{trial}}(\alpha)
        &= \mathrm{clip}\!\bigl(y_k + \alpha\, \Delta y,\; y^L, y^U\bigr),
        & \text{(scaled)} \\
    x_{\mathrm{trial}}(\alpha)
        &= \mathrm{clip}\!\bigl(x^{\mathrm{off}} + D\, y_{\mathrm{trial}}(\alpha),\; x^L, x^U\bigr),
        & \text{(physical)}
\end{align}
with $\alpha \in (0, 1]$. Because the transformation $y \mapsto x$ is
affine, walking along $\alpha$ traces the same physical line
$x_k + \alpha\,\Delta x$ regardless of which coordinate system one
parameterizes; the scaled parameterization is simply the bookkeeping
choice that makes the LP move limits unit-independent. A backtracking
line search starts from $\alpha = 1$ and successively shrinks
$\alpha \leftarrow \beta\,\alpha$ with $\beta \in (0, 1)$ (typically
$\beta = 0.5$) until an acceptance condition is met. The line search
is called \emph{inexact} because it does not minimize $\Phi$ exactly
along the direction; it only enforces a sufficient decrease, and that
decrease is measured on the \emph{physical} merit
$\Phi(x_{\mathrm{trial}}; M)$.

The implemented acceptance rule has two branches. The primary branch is
the Armijo-like sufficient-decrease condition
\citep[Sec.~3.1]{NocedalWright2006},\citep[Sec.~4.3]{MartinsNing2021}:
\begin{align}
    \Phi(x_{\mathrm{trial}}; M) \;\leq\;
    \Phi(x_k; M) \;-\; \eta\, \Delta\Phi_{\mathrm{pred}}(\alpha),
    \label{eq:armijo}
\end{align}
where $\eta \in (0, 1)$ is small (the implementation uses
$\eta = 10^{-4}$), and $\Delta\Phi_{\mathrm{pred}}(\alpha)$ is the merit
decrease predicted by the linearized model:
\begin{align}
    \Delta\Phi_{\mathrm{pred}}(\alpha)
    \;=\; \max\Bigl(0,\;
        \Phi(x_k; M) - \Phi^{\mathrm{lin}}_k(\alpha\Delta y; M)
    \Bigr),
\end{align}
with the linearized merit
$\Phi^{\mathrm{lin}}_k(\Delta y; M) =
 f_k + g_{y,k}^{\!\top}\Delta y + M \sum_i \max\bigl(0,
 -(c_{k,i} + (J_{y,k}\Delta y)_i)\bigr)$.

The secondary branch is a feasibility-progress condition. When the current
iterate is infeasible, the trial step is also accepted if it reduces the
total violation by at least a fraction $\gamma$:
\begin{align}
    V(x_{\mathrm{trial}}) \;\leq\; (1 - \gamma)\, V(x_k),
    \quad
    V(x) = \sum_i \max\bigl(0, -c_i(x)\bigr).
\end{align}
This rule comes from the filter and feasibility-restoration literature
\citep{FletcherLeyfferToint2002}: when the design is infeasible, a step
that meaningfully improves feasibility is useful even if the merit
function does not strictly satisfy~(\ref{eq:armijo}).

\subsubsection*{Why inexact line search makes the algorithm better}

\begin{itemize}[nosep,leftmargin=1.5em]
    \item \textbf{Globalization.} An LP step is built from a linear model
          that is only valid in a neighborhood of $x_k$. A full step can
          easily increase $f$ or worsen feasibility on the true
          nonlinear functions. Backtracking gives the solver a way to
          reduce the step until the nonlinear behavior catches up with
          the linear model.
    \item \textbf{Cheap.} Each backtrack step costs one nonlinear
          evaluation. No additional gradient is required, which is
          important here because gradients are obtained from finite
          differences and are therefore expensive (one finite-difference
          gradient costs $n$ evaluations of the structural model).
    \item \textbf{Convergence guarantees.} The Armijo-like rule is the
          mechanism by which classical line-search globalization theorems
          deliver global convergence to a stationary point under mild
          assumptions \citep[Thm.~3.2]{NocedalWright2006}. The same idea
          applied to a merit function is the standard globalization
          framework for SQP-type methods
          \citep[Sec.~18.5]{NocedalWright2006}.
    \item \textbf{Feasibility-progress branch.} Without it, the solver
          may reject useful restoration steps that temporarily increase
          $f$ but make the design feasible. This is the core idea of
          filter methods reused in a simpler form
          \citep{FletcherLeyfferToint2002}.
\end{itemize}

\subsection{Adaptive Move Limits --- A Trust-Region Strategy}
\label{sec:adaptive_move}

\subsubsection*{Theory}

The move limits $\delta$ act as a simple \emph{trust region}: they bound
how far in \emph{scaled} coordinates the LP is allowed to step. The
ratio that controls them, on the other hand, is built from
\emph{physical} merit values: this asymmetry is essential, because the
move limit must control LP conditioning (a scaled-space concern) while
$\rho$ must measure the agreement between linear model and true
nonlinear response (a physical-space concern). In trust-region theory
\citep[Ch.~6]{Conn2000},\citep[Ch.~4]{NocedalWright2006},
\citep[Sec.~4.5]{MartinsNing2021} the trust-region radius is updated by
comparing the actual reduction of the merit function to the reduction
predicted by the model, through the ratio
\begin{align}
    \rho_k
    \;=\;
    \frac{\Delta\Phi_{\mathrm{act}}}
         {\Delta\Phi_{\mathrm{pred}}}
    \;=\;
    \frac{\Phi(x_k; M) - \Phi(x_{k+1}; M)}
         {\Phi(x_k; M) - \Phi^{\mathrm{lin}}_k(\Delta y; M)}.
\end{align}
The standard interpretation is:
\begin{itemize}[nosep,leftmargin=1.5em]
    \item $\rho \approx 1$: the linear model agrees well with the
          nonlinear functions; the trust region can be enlarged.
    \item $0 < \rho \ll 1$: the model is over-optimistic; the trust
          region must be shrunk.
    \item $\rho \leq 0$: the step did not improve the merit at all;
          the step is rejected and the trust region is shrunk.
\end{itemize}

The implementation uses the classical update rule of
\citet[Algo.~4.1]{NocedalWright2006}, adapted to the SLP move limit:
\begin{align}
    \delta_{k+1}
    \;=\;
    \begin{cases}
        \min(a_{\mathrm{exp}}\, \delta_k,\; \delta_{\max}), & \rho > 0.75,\\
        a_{\mathrm{shr}}\, \delta_k,                        & \rho < 0.25,\\
        \delta_k,                                           & \text{otherwise},
    \end{cases}
\end{align}
with $a_{\mathrm{exp}} > 1$ and $0 < a_{\mathrm{shr}} < 1$. If the LP fails
or no backtracked trial point is acceptable, the move limit is reduced
unconditionally.

\subsubsection*{Why adaptive move limits are essential}

\begin{itemize}[nosep,leftmargin=1.5em]
    \item \textbf{Linear models are local.} The first-order model is
          accurate only in a neighborhood of $x_k$. The size of that
          neighborhood depends on the local curvature and on the
          finite-difference noise --- both of which the solver does not
          know in advance. The trust-region update infers this
          neighborhood from observed model quality.
    \item \textbf{Self-tuning.} Without adaptive move limits, the user
          would have to manually pick a step size that is small enough to
          be safe everywhere but large enough to converge in a reasonable
          number of iterations. The $\rho$-based update removes this
          manual tuning: the move limit shrinks automatically in
          difficult regions and grows automatically in well-behaved
          regions.
    \item \textbf{Robustness near nonlinear constraints.} Active
          structural constraints (buckling, stress) are highly nonlinear.
          When the LP repeatedly proposes steps that the line search
          rejects, the move limit shrinks and the LP is forced to operate
          in a regime where its linear model is reliable.
    \item \textbf{Theoretical foundation.} The $\rho$-based rule is part
          of the classical proof of trust-region global convergence
          \citep[Thm.~6.4.6]{Conn2000}. Although the present algorithm
          is not a pure trust-region method, the same logic underlies
          its globalization.
\end{itemize}

\subsection{Adaptive Penalty Update}
\label{sec:adaptive_penalty}

\subsubsection*{Theory}

The penalty $M$ controls how strongly the algorithm cares about feasibility
relative to the objective. In the slack-relaxed LP it discourages slack
usage. In the merit function it discourages stepping into infeasible
regions. The same parameter is used in both, which keeps the LP
subproblem and the merit acceptance criterion consistent.

A static value of $M$ is rarely satisfactory \citep[Sec.~17.3]{NocedalWright2006}:
\begin{itemize}[nosep,leftmargin=1.5em]
    \item too small: the optimizer prioritises objective reduction and
          accepts persistent infeasibility,
    \item too large: the LP is dominated by the slack penalty, the
          objective barely shapes the search direction, and
          ill-conditioning grows with $M$.
\end{itemize}
The classical adaptive rule increases $M$ whenever the current value
is insufficient to drive feasibility. The implementation uses the simple
monotone update
\begin{align}
    M_{k+1} \;=\; \min(a_M\, M_k,\, M_{\max}), \qquad a_M > 1,
\end{align}
triggered when either of the following holds after an iteration:
\begin{itemize}[nosep,leftmargin=1.5em]
    \item the LP needed non-negligible slack ($\sum s_i > \varepsilon_{\mathrm{slack}}$), or
    \item the iterate is still infeasible and the violation did not drop
          appreciably ($V(x_{k+1}) \geq 0.95\, V(x_k)$).
\end{itemize}

\subsubsection*{Why this is a good idea}

\begin{itemize}[nosep,leftmargin=1.5em]
    \item \textbf{Driving feasibility without manual tuning.}
          \citet[Sec.~5.4]{MartinsNing2021} argue that adaptive penalty
          methods are preferable to fixed penalties precisely because
          the user typically does not know what value of $M$ is large
          enough. The adaptive scheme grows $M$ until the LP and the
          merit function force the iterates toward feasibility.
    \item \textbf{Connection to exactness.} Increasing $M$ moves the
          merit function toward the regime in which it is an exact
          penalty function for the original problem
          \citep[Thm.~17.4]{NocedalWright2006}. In other words, the
          minimum of the merit function eventually coincides with the
          minimum of the constrained problem.
    \item \textbf{Diagnostic role of slacks.} Tying the penalty update to
          the LP slacks gives the solver a measurable, automatic trigger
          for ``the linearization cannot satisfy these constraints, so
          push harder''. This is far more reliable than a fixed schedule
          based on iteration count.
    \item \textbf{Bounded growth.} The cap $M_{\max}$ prevents
          ill-conditioning of the LP from running away in pathological
          cases.
\end{itemize}

\subsection{Convergence Criteria}
\label{sec:convergence}

The solver uses two complementary first-order convergence tests, plus a
graceful failure exit. Both tests require nonlinear feasibility ---
satisfied by the true model, not merely the linearized one --- so that a
collapsed move limit alone is never reported as success.

\paragraph{Primary criterion (small step at a feasible point).}
\begin{align}
    \max_i v_i(x_k) &\leq \varepsilon_{\mathrm{feas}}, \\
    \max_i s_i      &\leq \varepsilon_{\mathrm{slack}}, \\
    \|x_k - x_{k-1}\|_2 &\leq \varepsilon_x.
\end{align}
This is the operational meaning of ``the optimizer can no longer find a
useful improvement at a feasible iterate'' --- a standard step-norm test
\citep[Sec.~3.5]{MartinsNing2021}.

\paragraph{Secondary criterion (KKT-style first-order optimality).}
\begin{align}
    \max_i v_i(x_k) &\leq \varepsilon_{\mathrm{feas}}, \\
    \max_i s_i      &\leq \varepsilon_{\mathrm{slack}}, \\
    |f_k - f_{k-1}| &\leq \varepsilon_f, \\
    \|g_{y,k}\|_\infty &\leq \varepsilon_g.
\end{align}
This is the standard first-order necessary-optimality check
\citep[Sec.~12.3]{NocedalWright2006}: at a feasible local optimum, the
projected gradient must vanish. The implementation evaluates the gradient
norm in scaled space ($g_y$), which makes the tolerance $\varepsilon_g$
unit-independent.

\paragraph{Failure exit.}
If the move limit shrinks below $\delta_{\min}$ without the algorithm ever
satisfying either success criterion, the solver stops with a failure
message. This case means the globalization machinery could not find a
reliable step --- typically because the finite-difference gradient is
noisy, an active constraint is highly nonlinear, or the design is
trapped at a discontinuity. Importantly, the solver does \emph{not} report
this case as success, even though the iterates have stopped moving.

\subsubsection*{Why these criteria are appropriate}

\begin{itemize}[nosep,leftmargin=1.5em]
    \item \textbf{Feasibility is required.} Reporting an infeasible
          design as ``converged'' would be incorrect for a structural
          design problem, where every constraint corresponds to a
          physical requirement (load, deflection, buckling).
    \item \textbf{Two complementary checks.} The step-norm criterion
          catches the empirical situation in which the optimizer cannot
          move further. The first-order check captures the KKT-style
          definition of a local optimum. Either suffices to declare
          convergence; both must hold at the same feasible point in
          theory.
    \item \textbf{Scaled gradient.} Using $\|g_y\|_\infty$ instead of
          $\|g_x\|_\infty$ avoids declaring convergence merely because
          one variable has very small physical magnitude.
    \item \textbf{Honest failure.} Distinguishing ``stuck'' from
          ``optimal'' is essential for engineering decisions. The user
          can react to a ``stuck'' message by inspecting the recorded
          history (slacks, $\rho$, move limit), whereas a misleading
          ``optimal'' label could result in the wrong design being
          built.
\end{itemize}

\subsection{Algorithm Summary}

The complete iteration is summarized in Algorithm~\ref{alg:slp}. Every
ingredient described above appears as a clearly identifiable step.

\begin{table}[h]
\centering
\small
\begin{tabular}{p{0.04\textwidth} p{0.74\textwidth} p{0.12\textwidth}}
\toprule
\multicolumn{3}{l}{\textbf{Algorithm 1: Self-written SLP iteration}}\\
\midrule
\# & Step & Space \\
\midrule
1.  & Evaluate $f(x_k)$ and $c(x_k)$.                                                                  & physical \\
2.  & Estimate $g_x$, $J_x$ by bound-aware finite differences.                                          & physical \\
3.  & Transform gradients: $g_y = D g_x,\; J_y = J_x D$.                                                & scaled \\
4.  & Build and solve the slack-relaxed LP for $(\Delta y, s)$;
      decision variables and move limits in scaled space, RHS $c_k$
      and slacks $s$ in physical units.                                                                & mixed \\
5.  & Map step back: $\Delta x = D\,\Delta y$ for trial evaluations.                                    & physical \\
6.  & Compute the predicted merit decrease
      $\Delta\Phi_{\mathrm{pred}}(\alpha) = \Phi(x_k;M)
      - \Phi^{\mathrm{lin}}_k(\alpha\Delta y;M)$.                                                       & physical \\
7.  & Backtrack on $\alpha = 1, \beta, \beta^2, \dots$ until either
      Armijo merit decrease or feasibility-progress acceptance holds;
      tested on $\Phi(x_{\mathrm{trial}};M)$.                                                          & physical \\
8.  & If accepted, update $x_{k+1}, y_{k+1}, f_{k+1}, c_{k+1}$ and the
      history. Otherwise, mark the iteration rejected.                                                  & both \\
9.  & Compute $\rho_k=\Delta\Phi_{\mathrm{act}}/\Delta\Phi_{\mathrm{pred}}$
      and update $\delta$ in scaled space: enlarge if $\rho>0.75$,
      shrink if $\rho<0.25$.                                                                            & $\rho$: physical, $\delta$: scaled \\
10. & If slacks were active or feasibility stalled, increase $M$.                                       & physical \\
11. & Check $\|x_{k+1}-x_k\|_2\leq\varepsilon_x$ in physical space and
      $\|g_y\|_\infty\leq\varepsilon_g$ in scaled space at a feasible
      point; terminate with success or, if $\delta<\delta_{\min}$, with failure. & mixed \\
\bottomrule
\end{tabular}
\label{alg:slp}
\end{table}

\subsection{Default Parameter Values and Their Interpretation}

The solver exposes all numerical parameters as keyword arguments. The
defaults used in this thesis correspond to widely accepted choices in the
optimization literature \citep{NocedalWright2006,MartinsNing2021}.

\begin{table}[h]
\centering
\small
\begin{tabular}{lll}
\toprule
Parameter & Default & Interpretation \\
\midrule
$\delta_0 = \delta_{\max}$        & $0.10$    & initial / maximum scaled move limit\\
$\delta_{\min}$                   & $10^{-5}$ & minimum scaled move limit before failure\\
$a_{\mathrm{exp}}$                & $1.4$     & move-limit expansion factor\\
$a_{\mathrm{shr}}$                & $0.5$     & move-limit shrink factor\\
$M_0$                             & $10^{3}$  & initial merit / slack penalty\\
$a_M$                             & $10$      & penalty growth factor\\
$M_{\max}$                        & $10^{9}$  & penalty cap (avoids blow-up)\\
$\beta$                           & $0.5$     & backtracking shrink factor\\
$\eta$                            & $10^{-4}$ & Armijo sufficient-decrease constant\\
$\gamma$                          & $10^{-3}$ & feasibility-progress reduction fraction\\
$\varepsilon_{\mathrm{feas}}$     & $10^{-5}$ & nonlinear feasibility tolerance\\
$\varepsilon_{\mathrm{slack}}$    & $10^{-8}$ & LP slack tolerance\\
$\varepsilon_{x}$                 & $10^{-4}$ & step-size tolerance, physical $\|x_{k+1}-x_k\|_2$\\
$\varepsilon_{f}$                 & $10^{-3}$ & objective-change tolerance\\
$\varepsilon_{g}$                 & $10^{-6}$ & scaled-gradient tolerance\\
\bottomrule
\end{tabular}
\end{table}

\subsection{Connection to the Project}

In this thesis the solver is used to minimize the structural mass of a
support-structure cross-section subject to design-load and stability
inequality constraints. The choices made in the solver are motivated by
the project context as follows:

\begin{itemize}[nosep,leftmargin=1.5em]
    \item \textbf{Variable scaling} (Section~\ref{sec:scaling}) is used
          because the geometric variables of the cross-section span
          orders of magnitude. Without it, the LP would consistently
          favour the largest-magnitude variables irrespective of their
          actual influence on mass.
    \item \textbf{Slack-relaxed LP} (Section~\ref{sec:slacks}) is used
          because the initial design is generally infeasible: the
          designer typically starts from a guess that violates one or
          more constraints. The slacks allow the LP to make progress
          even from such a starting point.
    \item \textbf{Inexact backtracking line search}
          (Section~\ref{sec:line_search}) is used because finite-difference
          gradients are noisy: a full LP step routinely overshoots in the
          true nonlinear merit, and backtracking provides robustness at
          the cost of only a few extra structural evaluations per
          iteration.
    \item \textbf{Adaptive move limits} (Section~\ref{sec:adaptive_move})
          replace expensive manual tuning. The structural problem mixes
          well-behaved regions (e.g.\ where mass varies smoothly with
          radius) with strongly nonlinear regions (near buckling
          constraints), and a single fixed move limit cannot handle both.
    \item \textbf{Adaptive penalty}
          (Section~\ref{sec:adaptive_penalty}) ensures that final
          designs satisfy the structural constraints to engineering
          tolerance, not merely the linearized constraints.
    \item \textbf{Convergence criteria} (Section~\ref{sec:convergence})
          require true nonlinear feasibility before declaring success,
          which is a non-negotiable engineering requirement.
\end{itemize}

The recorded iteration history --- $\rho_k$, $\sum s_i$, $\delta_k$,
$M_k$, the backtracking step length $\alpha_k$ --- provides quantitative
feedback that is used in the results chapter to interpret the solver's
behavior on the structural problem and to compare it against
\textsc{PySLSQP}.

\subsection{Summary}

The implemented solver is a scaled, slack-relaxed Sequential Linear
Programming method with merit-function-based globalization, adaptive
move limits, and an adaptive penalty. Each ingredient corresponds to a
well-established result in the optimization literature
\citep{MartinsNing2021,NocedalWright2006,Conn2000,Fletcher1987,FletcherLeyfferToint2002}:
scaling for conditioning, slacks for LP feasibility, the $\ell_1$
merit function for acceptance, backtracking with Armijo for
globalization, the trust-region $\rho$-rule for adaptive move limits,
and an adaptive penalty for exactness. Together these mechanisms turn a
single linearization, which is only locally trustworthy, into a robust
nonlinear optimizer suitable for the structural design problem treated
in this thesis.

% =====================================================================
\begin{thebibliography}{9}

\bibitem{MartinsNing2021}
Martins, J.\,R.\,R.\,A.\ and Ning, A.,
\textit{Engineering Design Optimization}.
Cambridge University Press, 2021.

\bibitem{NocedalWright2006}
Nocedal, J.\ and Wright, S.\,J.,
\textit{Numerical Optimization}, 2nd ed.
Springer, 2006.

\bibitem{Fletcher1987}
Fletcher, R.,
\textit{Practical Methods of Optimization}, 2nd ed.
John Wiley \& Sons, 1987.

\bibitem{Conn2000}
Conn, A.\,R., Gould, N.\,I.\,M.\ and Toint, P.\,L.,
\textit{Trust-Region Methods}.
SIAM, 2000.

\bibitem{FletcherLeyfferToint2002}
Fletcher, R., Leyffer, S.\ and Toint, P.\,L.,
``On the global convergence of a filter-SQP algorithm,''
\textit{SIAM Journal on Optimization}, 13(1):44--59, 2002.

\end{thebibliography}

\end{document}
